\documentclass[10pt,a4paper]{amsart}
\usepackage[margin=19mm]{geometry}
\usepackage{amsmath,amssymb,mathtools}
\usepackage[scr=rsfs,cal=euler]{mathalfa}
\usepackage{xfrac}
\usepackage[unicode,hidelinks,bookmarks=false]{hyperref}
\newcommand{\ie}{i.e.\ }
\newcommand{\cf}{cf.\ }
\newcommand{\resp}{respectively\ }
\newcommand{\Subsection}{\subsection}
\providecommand{\nobreakdash}{\nobreak\hbox{-}}
\setlength{\emergencystretch}{2em}
\multlinegap0pt
\usepackage{amssymb}
\usepackage{mathtools}
\usepackage[shortlabels]{enumitem}
\usepackage{xcolor}
\usepackage{tikz}
\newtheorem{proposition}{Proposition}
\newcommand{\C}{\mathbb C}
\title{Log MMP Constraints on Curves with Cuspidal Singularities}
\author{Jenia Tevelev}
\date{Source: arXiv:2608.25031v1 (CC BY 4.0)}
\begin{document}
\maketitle

\noindent\emph{Source and licence.} Log MMP Constraints on Curves with Cuspidal Singularities. Version arXiv:2608.25031v1 (CC BY 4.0), distributed by arXiv under the Creative Commons Attribution 4.0 International licence, http://creativecommons.org/licenses/by/4.0/, as recorded in the arXiv licence metadata for this exact version and shown on the abstract page https://arxiv.org/abs/2608.25031v1. Full licence notice: \texttt{licenses/arxiv-2608.25031-CC-BY-4.0.txt}.\par\smallskip

\noindent\textit{Editorial scope.} Counted unit: the article's proposition prop:direct-sharpness (source0.tex lines 1881--1888). The article's complete original proof of it is delivered with it (source0.tex lines 1890--1981, one \texttt{proof} environment of 92 lines). No mathematics is written, repaired or re-derived: the statement and the proof are the article's own lines, in the article's own notation, and no retained line is substituted or re-flowed.  No further source-local context is needed: every cross-reference of every retained line resolves inside the two delivered spans.  Nothing was omitted from any retained span, so the package carries no omission record.  No retained line contains a citation, so the wrapper carries no bibliography and no bibliography entry.  No retained line contains a figure environment, an image-inclusion command or drawing code, so no image is delivered and the package declares no asset dependency.  Editorial formatting only, in addition: an AMS \texttt{amsart} preamble that restates the article's own declarations and macros used by the retained lines, namely the environments \texttt{proposition} on separate counters, together with the standard packages those lines need.  The retained environments print in this document as Proposition 1 in order of appearance, whereas the article prints the same items with the numbering of its own full document.  The complete native source of the article as distributed in the arXiv e-print for this exact version is bundled unchanged as \texttt{sources/208536/source0.tex}.
\begin{proposition}\label{prop:direct-sharpness}
For every coprime pair $1<p<q$, there is a smooth projective rational surface
$X_{p,q}$ containing a rational unicuspidal curve $C_{p,q}$ with Puiseux pair
$(p,q)$ and
\[
 C_{p,q}^{\,2}=m_{p,q}.
\]
\end{proposition}

\begin{proof}
Put
$s=p-p'$,
 $t=q-q'$,
 $k_p=\left\lfloor\frac ps\right\rfloor$,
 $k_q=\left\lfloor\frac qt\right\rfloor$,
$k=\max\{k_p,k_q\}$,
so that 
$m_{p,q}=pq+k$.

Draw the primitive vectors
\[
 v=(p,q),\qquad v_q=(p',t),\qquad v_p=(s,q'),
\]
which satisfy
\[
 \det(v_q,v)=\det(v,v_p)=1,
 \qquad v=v_q+v_p.
\]

Let $\sigma=\langle e_1,e_2\rangle$ be the first-quadrant cone.  We construct
a second regular cone $\tau$ containing $-v$ in its interior.
If $k=k_p$, set
$w_1:=k v_p-v$,
$w_2:=-v_p$.
Then
\[
 \det(w_1,w_2)=1,
 \qquad -v=w_1+k w_2.
\]
Moreover, the first coordinates of $w_1$ and $w_2$ are respectively
$ks-p\le0$ and $-s<0$.  Thus $\tau:=\langle w_1,w_2\rangle$ lies in the
closed left half-plane and has disjoint interior from $\sigma$.

If $k=k_q$, set instead
$
 w_1:=-v_q$,
 $w_2:=k v_q-v$.
Then
\[
 \det(w_1,w_2)=1,
 \qquad -v=k w_1+w_2,
\]
and the second coordinates of $w_1,w_2$ are $-t<0$ and $kt-q\le0$.
Hence $\tau:=\langle w_1,w_2\rangle$ lies in the closed lower half-plane and
again has disjoint interior from $\sigma$.  In the case of equality
$k_p=k_q$, either construction may be used.

Complete the two disjoint regular cones $\sigma$ and $\tau$ to a complete
regular fan by regularly subdividing the two remaining sectors.  Let
$X_{p,q}$ be the resulting smooth complete toric surface.

In the dense torus of $X_{p,q}$, consider the one-parameter subgroup
\[
 z\longmapsto (z^p,z^q),\qquad z\in\C^*,
\]
and let $C_{p,q}$ be its closure.  In the affine chart associated with
$\sigma$, the normalization is parametrized by
\[
 x=z^p,\qquad y=z^q,
\]
so the point $z=0$ is a cusp with one Puiseux pair $(p,q)$.  At $z=\infty$,
write $u=z^{-1}$.  Since the coordinates of $-v$ in the basis
$(w_1,w_2)$ are either $(1,k)$ or $(k,1)$, the local parametrization in the
chart of $\tau$ is
\[
 (u,u^k)\qquad\text{or}\qquad(u^k,u).
\]
Thus the curve is smooth at infinity.  On the torus it is an injective
immersion because $\gcd(p,q)=1$.  Consequently $C_{p,q}$ is rational and
has no singularity other than the prescribed cusp.

It remains to compute its square.  The cusp has
$\delta=(p-1)(q-1)/2$, hence
\[
 2p_a(C_{p,q})-2=pq-p-q-1.
\]
The curve meets the toric boundary only at $0$ and $\infty$.  Its local
intersection multiplicities with the two boundary components at $0$ are
$p,q$, and those at $\infty$ are $1,k$.  Since
$-K_{X_{p,q}}$ is the sum of the torus-invariant boundary divisors,
\[
 -K_{X_{p,q}}\cdot C_{p,q}=p+q+k+1.
\]
Adjunction therefore gives
$
 C_{p,q}^{\,2}
 =pq-p-q-1+p+q+k+1
 =pq+k
 =m_{p,q}.
$
\end{proof}

\end{document}
